\documentclass[11pt]{article}

% --------------------------------------------------
% Packages
% --------------------------------------------------
\usepackage{geometry}
\usepackage{setspace}
\usepackage{amsmath,amssymb,amsthm}
\usepackage{mathtools}
\usepackage{hyperref}
\usepackage{csquotes}
\usepackage{microtype}

\geometry{margin=1in}
\setstretch{1.15}

% --------------------------------------------------
% Theorem Environments
% --------------------------------------------------
\newtheorem{definition}{Definition}[section]
\newtheorem{axiom}{Axiom}[section]
\newtheorem{proposition}{Proposition}[section]
\newtheorem{theorem}{Theorem}[section]
\newtheorem{lemma}{Lemma}[section]
\newtheorem{corollary}{Corollary}[section]

\title{Replay and Invariance:\\
Toward a Structural Geometry of Semantic Systems}
\author{Flyxion}
\date{\today}

\begin{document}
\maketitle

\begin{abstract}
Modern physics and mathematics are unified not by ontology but by invariance. Physical laws, geometric structures, and mathematical objects are defined by what remains unchanged under admissible transformations. This paper develops a parallel structural framework for semantic systems, grounded in the notion of replayable construction histories. We treat replay as an axiomatic constraint analogous to path-independence or flatness conditions in geometry, and abstraction as the formation of equivalence classes invariant under replay. Without appealing to representational semantics or physical interpretation, we show that semantic stability admits a geometric description in terms of quotient structures, invariants, and structural phase transitions. The resulting framework positions meaning as an emergent invariant of construction, aligning semantic systems with the structural core of modern mathematics and physics.
\end{abstract}

% ==================================================
\section{Introduction}
% ==================================================

Across physics and mathematics, explanatory power arises not from enumerating objects but from identifying invariants. From classical mechanics to gauge theory, from topology to algebra, what distinguishes structure from accident is invariance under transformation. Meaningful structure persists; noise does not.

Semantic systems—whether scientific theories, mathematical formalisms, or computational constructions—exhibit a similar distinction. Some differences matter; others do not. Yet unlike physical systems, semantics has traditionally been treated as interpretive or representational rather than structural.

This paper advances a different position: semantic structure can be treated geometrically, as an invariant arising from replayable construction histories. Replay is not a metaphor for time reversal or simulation; it is an axiomatic constraint ensuring that meaning is defined by necessity rather than convenience. Abstraction, on this view, is not compression or symbol manipulation, but the formation of equivalence classes invariant under replay.

The goal is not to reduce semantics to physics, but to show that both obey the same structural logic.

% ==================================================
\section{Structural Invariance as a Unifying Principle}
% ==================================================

In mathematics, an object is defined by its invariants under a group of transformations. In physics, laws are constraints invariant under changes of frame, gauge, or coordinates. In both cases, structure is identified with stability under admissible transformations.

\begin{definition}
An invariant is a property preserved under a specified class of transformations.
\end{definition}

The choice of admissible transformations is not arbitrary; it encodes the theory’s commitments. Euclidean geometry permits rigid motions; gauge theories permit local phase shifts; topology permits continuous deformation.

Semantic systems likewise admit transformations: refactorings, re-derivations, alternative constructions. The question is which of these transformations preserve meaning.

% ==================================================
\section{Replay as an Axiomatic Constraint}
% ==================================================

We introduce replay not as a dynamical process but as a structural axiom.

\begin{axiom}
All admissible semantic transformations must preserve replayability of construction histories.
\end{axiom}

Replayability asserts that a construction can be re-instantiated from its generative history without ambiguity. This axiom plays a role analogous to integrability conditions in differential geometry or path-independence in conservative systems.

\begin{proposition}
Replayability enforces a notion of semantic flatness: equivalent constructions yield identical consequences under all admissible replays.
\end{proposition}

This flatness does not imply triviality. Just as flat connections may admit rich global structure, replayable systems may support deep abstraction.

% ==================================================
\section{Equivalence Classes and Quotient Structure}
% ==================================================

Abstraction arises when distinct constructions are identified as semantically equivalent. Structurally, this corresponds to forming a quotient space.

\begin{definition}
A semantic equivalence class is a set of construction histories whose replayed consequences are identical under all admissible perturbations.
\end{definition}

The space of constructions is partitioned into equivalence classes. The abstract object is not any representative construction, but the class itself.

\begin{proposition}
Semantic quotienting reduces degrees of freedom while preserving invariant structure.
\end{proposition}

This mirrors familiar constructions in mathematics, where quotienting by symmetry removes redundancy while revealing essential form.

\begin{corollary}
Abstraction is neither lossless nor arbitrary: it is constrained by the requirement of replay invariance.
\end{corollary}

% ==================================================
\section{Geometry of Semantic Space}
% ==================================================

Although semantic spaces need not be manifolds, they admit geometric intuition. Distance may be defined operationally as divergence under replay; curvature as sensitivity to perturbation; stability as resistance to equivalence class fragmentation.

\begin{definition}
Semantic distance between two constructions is the minimal perturbation under which their replayed consequences diverge.
\end{definition}

Regions of low curvature correspond to robust abstractions; regions of high curvature indicate unstable or context-dependent meaning.

This geometry is not imposed but induced by replay constraints.

% ==================================================
\section{Phase Transitions and Structural Change}
% ==================================================

In many systems, abstraction appears suddenly rather than gradually. Structurally, this corresponds to a phase transition in equivalence structure.

\begin{proposition}
A semantic phase transition occurs when the topology of equivalence classes changes discontinuously under continuous variation of construction parameters.
\end{proposition}

Examples include the emergence of conserved quantities, the stabilization of mathematical formalisms, or the crystallization of scientific theories.

Such transitions are structural, not psychological or sociological.

% ==================================================
\section{Relation to Physical Theories}
% ==================================================

The framework developed here is deliberately interpretation-agnostic. Nevertheless, it resonates with several physical notions.

Replay resembles reversibility constraints but does not require time symmetry. Semantic flatness parallels gauge invariance without invoking fields. Quotienting echoes symmetry reduction without presupposing physical symmetry groups.

\begin{proposition}
Semantic invariance is orthogonal to physical ontology: it constrains structure, not substance.
\end{proposition}

This allows semantic systems to be compared structurally to physical systems without reductionism.

% ==================================================
\section{Mathematics as Stabilized Replay}
% ==================================================

Mathematical objects may be viewed as maximally stable semantic equivalence classes. Proofs are constructions; theorems are invariants under replay.

\begin{proposition}
A mathematical object is semantically defined by the invariance of its consequences under replay of admissible derivations.
\end{proposition}

This explains why mathematical meaning is robust to notation, presentation, and proof strategy.

% ==================================================
\section{Limitations and Open Directions}
% ==================================================

This paper does not provide a full formalization of semantic geometry. Nor does it claim that all semantic systems admit smooth structure.

Rather, it establishes replay invariance as a minimal axiom from which geometric reasoning becomes possible.

Future work may explore categorical formulations, topological invariants, or connections to physical information constraints, but such developments are extensions, not prerequisites.

% ==================================================
\section{Conclusion}
% ==================================================

By treating replay as an axiomatic constraint and abstraction as invariant quotienting, semantic systems can be understood within the same structural framework that underlies modern physics and mathematics. Meaning emerges not from representation or interpretation, but from stability under transformation.

This perspective does not reduce semantics to geometry or physics. It shows that all three share a deeper commitment: structure is what survives change.

\appendix

\section{Replay Axiomatics}

We begin by formalizing replay as a structural constraint on construction systems. The aim is to isolate the minimal assumptions required to speak meaningfully about abstraction, equivalence, and invariance, without presupposing any particular computational, physical, or representational substrate.

\begin{definition}[Construction System]
A construction system is a tuple
\[
\mathcal{C} = (\mathcal{E}, \mathcal{H}, \mathcal{R})
\]
where $\mathcal{E}$ is a set of primitive events, $\mathcal{H}$ is the set of finite event histories (finite sequences over $\mathcal{E}$), and $\mathcal{R} : \mathcal{H} \to \mathcal{S}$ is a replay map assigning each history a semantic outcome in a space $\mathcal{S}$.
\end{definition}

The replay map need not be injective. Multiple distinct histories may correspond to the same semantic outcome.

\begin{axiom}[Replay Determinacy]
For any history $h \in \mathcal{H}$, the replay map $\mathcal{R}(h)$ is uniquely determined.
\end{axiom}

This axiom excludes stochastic or underdetermined replay, but does not constrain how histories are generated, nor does it require that histories be physically reproducible.

\begin{axiom}[Replay Closure]
If $h \in \mathcal{H}$, then any prefix $h' \preceq h$ is also an element of $\mathcal{H}$.
\end{axiom}

Replay closure ensures that semantic construction is well-founded and that intermediate stages of construction are themselves meaningful within the system.

\section{Semantic Equivalence and Quotient Structure}

Abstraction is formalized as a quotient induced by replay invariance.

\begin{definition}[Replay Equivalence]
Two histories $h_1, h_2 \in \mathcal{H}$ are replay-equivalent, written $h_1 \sim h_2$, if and only if
\[
\mathcal{R}(h_1) = \mathcal{R}(h_2).
\]
\end{definition}

\begin{proposition}
Replay equivalence is an equivalence relation on $\mathcal{H}$.
\end{proposition}

\begin{proof}
Reflexivity and symmetry follow immediately from equality. Transitivity follows from the transitivity of equality in the semantic space $\mathcal{S}$.
\end{proof}

\begin{definition}[Semantic Quotient Space]
The semantic quotient space is defined as
\[
\mathcal{Q} = \mathcal{H} / \sim,
\]
whose elements are equivalence classes of replay-equivalent histories.
\end{definition}

The quotient is induced by necessity, not by optimization or representational convenience.

\begin{proposition}
The canonical projection $\pi : \mathcal{H} \to \mathcal{Q}$ is surjective but generally not injective.
\end{proposition}

This non-injectivity captures the core notion of abstraction: many constructions collapse onto a single semantic representative.

\section{Replay Flatness and Path Independence}

Replay induces a flatness condition analogous to path independence in geometric and algebraic settings.

\begin{definition}[Semantic Flatness]
A construction system is semantically flat if, for any two histories $h_1, h_2 \in \mathcal{H}$ that differ only by permutation of independent events, the replay outcomes coincide:
\[
\mathcal{R}(h_1) = \mathcal{R}(h_2).
\]
\end{definition}

\begin{proposition}
Semantic flatness implies that semantic outcome depends only on equivalence class, not on construction path.
\end{proposition}

\begin{corollary}
Violations of semantic flatness correspond to semantic curvature, in which local construction order affects global outcome.
\end{corollary}

No smoothness, differentiability, or manifold structure is assumed. The analogy to curvature is purely structural.

\section{Semantic Distance and Stability}

To reason about robustness and invariance, we introduce a minimal notion of semantic distance.

\begin{definition}[Perturbation]
A perturbation is a function
\[
\delta : \mathcal{H} \to \mathcal{H}
\]
that modifies a history by a bounded change, such as event substitution, deletion, or insertion.
\end{definition}

\begin{definition}[Semantic Distance]
The semantic distance between histories $h_1$ and $h_2$ is defined as
\[
d(h_1, h_2) = \inf \{ \|\delta\| \mid \mathcal{R}(\delta(h_1)) \neq \mathcal{R}(\delta(h_2)) \}.
\]
\end{definition}

This function defines a pseudometric on equivalence classes; it need not satisfy the triangle inequality globally.

\begin{definition}[Semantic Stability]
An equivalence class $[h] \in \mathcal{Q}$ is semantically stable if there exists $\epsilon > 0$ such that for all $h' \notin [h]$,
\[
d(h, h') > \epsilon.
\]
\end{definition}

Stable equivalence classes correspond to abstractions that persist under bounded perturbations.

\section{Structural Phase Transitions}

We now formalize abstraction emergence as a structural phase transition.

\begin{definition}[Control Parameter]
A control parameter $\lambda$ is a continuous parameter governing admissible perturbations or construction constraints within a family of construction systems $\mathcal{C}_\lambda$.
\end{definition}

\begin{definition}[Semantic Phase Transition]
A semantic phase transition occurs at $\lambda = \lambda_c$ if the topology or connectivity structure of the quotient space $\mathcal{Q}_\lambda$ changes discontinuously at $\lambda_c$.
\end{definition}

\begin{proposition}
Semantic phase transitions correspond to qualitative changes in equivalence-class stability, connectivity, or degeneracy.
\end{proposition}

This definition captures abstraction emergence without appeal to thermodynamic entropy or statistical mechanics.

\section{Mathematical Objects as Invariants}

Mathematical objects arise naturally as maximally stable semantic invariants.

\begin{definition}[Admissible Derivation]
An admissible derivation is a history whose events correspond to rule-governed inferential steps within a formal system.
\end{definition}

\begin{proposition}
A mathematical object corresponds to a semantic equivalence class invariant under all admissible derivations.
\end{proposition}

\begin{corollary}
Mathematical meaning is independent of representational form, proof path, or syntactic encoding.
\end{corollary}

This provides a structural account of mathematical objecthood without epistemic or psychological assumptions.

\section{Scope and Non-Commitments}

To avoid misinterpretation, we make the following clarifications explicit.

No assumption is made that semantic spaces are smooth manifolds. No reliance is placed on category theory, though the framework is compatible with later categorical refinement. No identification is made between semantic space and physical spacetime. No claim is made that replay is physically realizable in all systems. No reduction of meaning to entropy or information is assumed.

\begin{axiom}[Minimality]
Only structural commitments explicitly stated in this paper are assumed.
\end{axiom}

\end{document}

